\section{Scope, conventions, and the precise advances}
\label{sec:scope}

The aim is to turn several questions already present in the ProveIt
research stream into exact theorems. The most direct antecedents are
the mixed spectral report, the independent-order Hurwitz report, and
the raw harmonic-power section of the exact resonant report
\cite{RepoMixed,RepoIndependent,RepoExactResonant}. These sources
already contain substantial analytic continuation, Gamma normalization,
and finite-part machinery. We use it with explicit attribution, supply
self-contained proofs of the identities needed here, and resolve
specific questions about singularities, spans, and collision constants.

The literature background is classical: Gamma and polygamma
recurrences, Hurwitz-zeta continuation, Dougall's summation, and
polylogarithmic Mellin transforms
\cite{DLMFPolygamma,DLMFHurwitz,DLMFDougall,DLMFPolylog}.
The Arakawa--Kaneko functions are established objects
\cite{ArakawaKaneko,KanekoTsumura}. Our claims of an advance are
relative to the inspected ProveIt material. We do not claim that every
special-value consequence is absent from the wider literature.

\subsection{What is established}

\begin{longtable}{@{}p{0.20\textwidth}p{0.36\textwidth}p{0.38\textwidth}@{}}
\toprule
Problem & Antecedent & Result of this continuation\\
\midrule
\endhead
Shift singularities & A depth-dependent continuation half-plane;
the broader relative ordered Hurwitz character and its transport.
& Exact local amplitudes, boundary constants, sharp Taylor radii,
normalized primitives, and the finite-to-infinite functional rank
transition; Section~\ref{hsh:section}.\\[6pt]
Polygamma moments & A two-parameter Dougall identity and the request
to separate moments at weights six and eight.
& The full mixed-Bell rank in every weight, exact elimination
matrices, the three-dimensional summation image, and its full kernel;
Section~\ref{sec:bell}.\\[6pt]
Centered harmonic powers & Regularity of $F_p(-2r,1/2)$ and a
multiple-zeta expression obtained by finite word calculations.
& A triangular generator for every $r$ and all $p$, and one finite
Euler-sum space independent of $r$ at fixed $p$;
Section~\ref{hpnew:section}.\\[6pt]
Geometric collisions & Fixed-rate triple digamma formulas and
questions about nested rates and higher argument derivatives.
& An exact all-factor polygamma correction with jointly holomorphic
remainder, oriented harmonic hierarchy laws, and explicit derivative
examples; Section~\ref{hc:section}.\\
\bottomrule
\end{longtable}

The harmonic generator answers the all-power minus-two target and
gives every further even moment by finitely many triangular steps.
It is not a single closed two-variable expression in both the power
and the moment order. A companion result proves regularity for
polynomials in odd generalized harmonic orders and evaluates each
single odd-order family in a fixed finite space of even zeta values.
The collision theorem treats Stieltjes index
zero with arbitrary argument derivatives. Positive Stieltjes indices,
which introduce logarithmic local singularities, remain a separate
extension. The Bell rank concerns the stated formal polynomial rows;
it does not decide arithmetic independence of special values.
These distinctions delimit what has actually been proved.

\subsection{Functions and derivative conventions}

We use the ordinary Gamma function, $\psi=\Gamma'/\Gamma$, and
$\psi^{(q)}=\partial_a^q\psi$. The generalized harmonic numbers are
\[
 H_n^{(q)}(a)=\sum_{j=0}^{n-1}(a+j)^{-q},\qquad
 H_n^{(q)}=H_n^{(q)}(1),\qquad H_n=H_n^{(1)}.
\]
For an integer $q\ge1$ they satisfy
\begin{equation}\label{scope:harmonic-polygamma}
 H_n^{(q)}(a)=\frac{(-1)^{q-1}}{(q-1)!}
       \bigl\{\psi^{(q-1)}(a+n)-\psi^{(q-1)}(a)\bigr\}.
\end{equation}
The same formulas extend in $a$ away from their poles, with branches
specified when nonintegral powers occur. Relation
\eqref{scope:harmonic-polygamma} follows by telescoping the polygamma
recurrence \cite{DLMFPolygamma}.

For $\Re s>1$ and $a>0$,
$\zeta(s,a)=\sum_{n\ge0}(n+a)^{-s}$, and
$\zeta(s)=\zeta(s,1)$. Generalized Stieltjes constants have the
fixed convention
\begin{equation}\label{scope:Stieltjes}
 \zeta(1+v,a)=\frac1v+
       \sum_{j\ge0}\frac{(-1)^j}{j!}\gamma_j(a)v^j,
 \qquad \gamma_0(a)=-\psi(a),\qquad\gamma_j=\gamma_j(1).
\end{equation}
Thus $\gamma_j^{(q)}(a)$ always denotes an argument derivative,
whereas $\partial_s^k\zeta(s,a)$ is a spectral derivative.
The index $j$ in $\gamma_j$ is discrete and is never differentiated.
Similarly, a formula known at integer spectral labels $s=-m$ cannot
be differentiated with respect to $m$ to infer a spectral derivative.

Our multiple-polylogarithm convention is
\[
 \Li_{k_1,\ldots,k_d}(z)=
 \sum_{n_1>\cdots>n_d\ge1}
       \frac{z^{n_1}}{n_1^{k_1}\cdots n_d^{k_d}}.
\]
For admissible indices, its value at $z=1$ is
$\zeta(k_1,\ldots,k_d)$; weight means $k_1+\cdots+k_d$.
The one-index case is the ordinary polylogarithm
\cite{DLMFPolylog}. The Lerch transcendent is initially
\[
 \Phi(z,s,a)=\sum_{n\ge0}\frac{z^n}{(n+a)^s},\qquad |z|<1,
\]
with analytic continuation where required
\cite{DLMFLerch}. In particular $z\Phi(z,s,1)=\Li_s(z)$.

\subsection{Three different meanings of a displayed infinite expression}

Most sums in Section~\ref{sec:bell}, including the tail-square
identity, converge absolutely in the ordinary sense. They need no
regularization. In contrast,
\[
 F_p(-2r,1/2)
\]
means the value of the meromorphic continuation of
$\sum_{n\ge0}H_n^p(n+1/2)^{-s}$ at the regular point $s=-2r$.
For $p\ge1$ this is not the sum of an ordinarily convergent numerical
series after direct substitution of $s=-2r$.

For a function with a local Laurent--log expansion, $\CT_z$ extracts
the coefficient of $z^0(\log z)^0$. A Hadamard finite-part integral
is the constant term in its cutoff expansion with the cutoff coordinate
held fixed. At distinct interior seams we use independent ambient
coordinates $x-x_i$ and take their joint constant term before a
collision parameter tends to zero. The subsequent $\CT_\epsilon$
in Section~\ref{hc:section} is a second operation, on the resulting
function of the geometric shifts. Changing either coordinate can
change this constant. The exact correction formulas record that
dependence.

Branches are also part of the statements. A principal shift
continuation, a continuation around a negative integer, and a real
collision path need not have the same analytic germ. All local powers
are written using an explicitly chosen logarithm. Fixed-endpoint
primitives are normalized to vanish at their stated base point, so
an integration constant cannot absorb a resonance discrepancy.

\subsection{Proof strategy and verification}

Each family is handled by first retaining a full analytic parameter
family, then making the desired specialization. Finite transport
isolates a shift singularity; a Gamma quotient generates the Bell
moments; an exponent parameter linearizes raw harmonic powers; and
partial fractions isolate the singular collision pieces. Normal
convergence or a holomorphic remainder is proved before extracting
coefficients. This order of operations is essential at resonances.

Exact symbolic checks replay coefficient identities and rational
linear algebra. Independent numerical evaluations start from different
representations and diagnose signs, constants, branches, and omitted
terms. They supplement the proofs. The numerical tails and quadratures
are not interval certificates, and their small residuals are not used
to justify an infinite identity. Section~\ref{sec:audit} records the
actual checks and the boundaries of the source audit.
